% No preamble. Prefix cmt-. True infinite Maxwell coefficient; not finite cutoff beta.
\section{Exact classical preservation of the Maxwell coefficient}
\label{cmt-section}

This section concerns the invariant infinite interaction quotient
on one selected joint classical background system. The coefficient
\(\beta_M\) is the true infinite total-kernel Maxwell moment.
It is not the finite symmetric-cutoff ratio \(\beta_{M,N}\).
Fix the primitive ratio, grouping \(J\), activity exponent \(a_0\),
and one positive field radius \(b\). The actual homogeneous
reference and projection satisfy
\begin{equation}
 \beta_M(S_M)=1,\qquad \|S_M\|_b\le C_Sb^2,\qquad
 P_M=I-S_M\beta_M,\qquad \|P_M\|\le C_{\rm ext}(J).
 \label{cmt-reference}
\end{equation}
The actual pre-gauge substitution has the scalar identities
\begin{equation}
 \mathcal T_{M,L}S_M=S_{ML},\qquad
 \mathcal T_{MR,L}\mathcal T_{M,R}
                       =\mathcal T_{M,RL}.
 \label{cmt-semigroup}
\end{equation}
They are identities of normalized interactions, not of chosen
support coefficient arrays.

The complete-field pre-gauge family obeys, for every fixed
prescribed finite classical reserve \(A\),
\begin{equation}
 \sup_i\sum_X e^{\kappa\tau(\{i\}\cup X)+A|X|}
       \|a_{M,L;i,X}\|_{C^{1,\alpha};b}
           \le C_* b/L,\qquad 0<b\le r_A,
 \label{cmt-uniform-family}
\end{equation}
on one outer radius \(r_A>0\) independent of \(M,L\).
The entire output fiber and its physical first-gradient
H\"older quotients are inside the support sum, and the
coefficients agree on overlapping collars. Fix \(J\) and
the reserves so that \(2(a_0+1)/J\le\kappa\) and
\(A\ge4(a_0+1)\).

The argument uses countably many arbitrarily large endpoint
powers. They can be included in the chosen joint system by
a further diagonal extraction with the same \(r_A\);
previously fixed limits remain unchanged. The exact finite
comparison and controlled gauge identities then give
\eqref{cmt-semigroup} on this enlarged joint system.
No independence from unrelated subsequences is asserted.

\subsection{The quadratic estimate needs no large field-scaling disk}
\label{cmt-quadratic-section}

Let \(y=P_Mx\) for an invariant class \(x\). Global semisimple
color invariance gives \(D\Phi_y(0)=0\). The exact infinite Ward
classification and hypercubic symmetry give
\begin{equation}
 \Gamma_2^y(p)=O(|p|^4),\qquad
 \|K_2(y)\|_{\sigma(J),{\rm row}}
                         \le C(J)b^{-2}\|y\|_b.
 \label{cmt-kernel}
\end{equation}
The fixed midpoint phases preserve the order of vanishing in
tail coordinates. Define \(Q_y(A)=K_2(y)[A,A]/2\).
For compact coarse \(w\), the pre-gauge image is exponentially
decaying and summable, so \(Q_y(a_{M,L}(w))\) is well defined.

The quadratic part of the spatial localization proof yields
a connected family \(Q_{y,L}\) representing that entire
substituted quadratic polynomial, with
\begin{equation}
 \|Q_{y,L}\|_{a_0,b}
       \le C_2(J)L^{-\alpha}\|y\|_b,
       \qquad L\ge L_0(J),\quad 0<b\le r_A.
 \label{cmt-quadratic-bound}
\end{equation}
Indeed the divided-difference factorization is
\(K_2=d^-\widetilde H d^+\), with
\(\sum_r\widetilde H(r)=0\), and gives the exact expression
\[
 \frac12\sum_{x,r,\rho,\tau}
 [D_\rho a(x)]^T\widetilde H_{\rho\tau}(r)
                    [D_\tau a(x+r)-D_\tau a(x)].
\]
The complete family bounds supply \(CbL^{-2}\) for the first
difference and \(CbL^{-2-\alpha}|r|^\alpha\) for the nearby
gradient difference. The exponential kernel tail controls
the far terms. The \((JL)^4\) anchor count leaves
\(L^{-\alpha}\). All constants are fixed after \(J\) and
before \(L\).

Only \(a_{M,L}(w)\) on the original \(b\)-ball is used.
There is no substitution \(a_{M,L}(\lambda w)\) with
\(|\lambda|\) proportional to \(L\). Thus the high-jet
condition \(bL\le\text{constant}\) is absent from
\eqref{cmt-quadratic-bound}.

Although \(Q_{y,L}\) need not be invariant as a whole nonlinear
functional, the chosen Hessian moment extends to every analytic
interaction quotient as a bounded linear functional:
\begin{equation}
 |b_{\rm mom}(Z)|\le C_\beta(J)b^{-2}\|Z\|_b.
 \label{cmt-moment-bound}
\end{equation}
On invariant classes it is \(\beta\). This norm bound is just
a fixed second derivative and a fixed second spatial moment;
it does not require a Ward identity for \(Z\).

\subsection{The transported Hessian and the semigroup limit}
\label{cmt-preservation-section}

Since \(a_{M,L}(0)=0\) and \(D\Phi_y(0)=0\), the chain rule gives
\begin{equation}
 D^2(\Phi_y\circ a_{M,L})(0)[u,v]
   =D^2\Phi_y(0)[Da_{M,L}(0)u,Da_{M,L}(0)v].
 \label{cmt-chain-rule}
\end{equation}
The right side is also the Hessian of
\(Q_y\circ a_{M,L}\). Higher input field degrees cannot
contribute to degree two, and the possible
\(D\Phi_y(0)D^2a_{M,L}(0)\) term is zero.
Differentiation is justified first on finitely supported
directions by local normal convergence, then for the moment
by the anchored exponential kernel bounds. Hence
\begin{equation}
 \beta_{ML}(\mathcal T_{M,L}y)=b_{\rm mom}(Q_{y,L}).
 \label{cmt-hessian-identification}
\end{equation}
Combining \eqref{cmt-quadratic-bound} and
\eqref{cmt-moment-bound}, followed by the exact reference
transport, proves
\begin{equation}
 |\beta_{ML}(\mathcal T_{M,L}x)-\beta_M(x)|
       \le C_m(J)b^{-2}L^{-\alpha}\|x\|_b.
 \label{cmt-asymptotic}
\end{equation}
This estimate is uniform in preceding depth, at fixed \(b>0\).

\begin{theorem}[Exact infinite Maxwell preservation]
\label{cmt-theorem}
For every fixed admissible ratio \(R\) and every invariant
class \(x\) on this common domain,
\begin{equation}
       \beta_{MR}(\mathcal T_{M,R}x)=\beta_M(x).
 \label{cmt-exact}
\end{equation}
\end{theorem}
\begin{proof}
Put \(z=\mathcal T_{M,R}x\). The fixed-\(R\) classical
transport is bounded, say
\(\|z\|_b\le B_R(J)\|x\|_b\). Apply
\eqref{cmt-asymptotic} to \(z\) at depth \(MR\) and
endpoint \(L\), and to \(x\) at depth \(M\) and endpoint
\(RL\). By \eqref{cmt-semigroup}, their leading transported
moments agree. Therefore
\begin{equation}
 |\beta_{MR}(\mathcal T_{M,R}x)-\beta_M(x)|
 \le C_m(J)b^{-2}
       \bigl[B_R(J)L^{-\alpha}+(RL)^{-\alpha}\bigr]\|x\|_b.
 \label{cmt-comparison}
\end{equation}
The left side is independent of \(L\). Let \(L\) tend through
the admissible endpoint powers to infinity in the joint
system.
\end{proof}

The proof needs no convergence of the full interactions
\(\mathcal T_{M,L}x\) as \(L\to\infty\), and no bound on
their whole norms uniform in \(L\). Only the fixed-\(R\)
map, the scalar semigroup, and the quadratic-only bound
are used.

There is no hidden domain shrink. The actual pre-gauge
holomorphy and \eqref{cmt-uniform-family} have a fixed
primitive-scale radius. A real minimum comparison at
one large endpoint can initially require a smaller real
germ because physical norms rescale. Both scalar sides
of \eqref{cmt-semigroup}, however, are holomorphic on
their specified common \(b\)-domain; equality on that
nonempty germ extends throughout the connected domain.
On compact field subspaces the anchored and summable-field
bounds justify that continuation. The high-jet dilation
radius plays no role. Countably many endpoints can
therefore be included without replacing \(b\) by a
sequence tending to zero.

\subsection{The original Maxwell coordinate isolates the leading drift}
\label{cmt-drift-section}

Equations \eqref{cmt-exact} and \eqref{cmt-semigroup} imply
\begin{equation}
 \beta_{MR}\mathcal T_{M,R}=\beta_M,\qquad
 \mathcal T_{M,R}P_M=P_{MR}\mathcal T_{M,R}.
 \label{cmt-commutation}
\end{equation}
Thus the classical leakage on \(\ker\beta_M\) is exactly
zero for this infinite system. An adapted-coordinate
correction introduced solely to remove that leakage
can be chosen zero; the general coordinate theorem
remains valid for other operator families.

For the fixed-coupling quantum map
\(\mathcal F_{M,g,T}=\mathcal T_{M,R}+\mathcal G_{M,g,T}\),
let \(x=P_Mx\), \(\|x\|_b\le C_{\rm int}b^2\). The parity
theorem defines
\(D_M=\Delta_M+\mathcal Q_{{\rm det},M}\), as in
\eqref{cgp-leading-definition}, and gives the inner bound
\eqref{cgp-inner-bound}. The input Lipschitz estimate gives
\[
 \|\mathcal G_{M,g,T}(x)-\mathcal G_{M,g,T}(0)\|_b
      \le C_{\rm quant}(gT/b)C_{\rm int}b^2.
\]
Applying the original Maxwell moment, with
\(\|\beta_{MR}\|\le C_\beta(J)b^{-2}\), yields
\begin{align}
 &|\beta_{MR}(\mathcal F_{M,g,T}(x))
                         -\beta_{MR}(D_M)|\notag\\
 &\quad\le
   \frac{2C_\beta C_{\rm par}}{\rho^2}
                         [\tau(T)+g^2T^6]
       +C_\beta C_{\rm quant}C_{\rm int}\,gT/b.
 \label{cmt-drift-bound}
\end{align}
All constants are uniform in preceding depth on the fixed
classical system. The field radius \(b\) is fixed before
the real threshold \(T=T(g)\) and \(g\to0\) are considered.
For an actual density update
\(\gamma_{\rm next}=\gamma-\beta_{MR}(\mathcal F(x))\),
the leading decrement in the original coupling coordinate
is therefore \(\beta_{MR}(D_M)\), with the displayed
vanishing error. No adapted marginal coordinate is needed
for this conclusion.

\subsection{Finite and physical limitations}
\label{cmt-scope}

The finite coefficient \(\beta_{M,N}\) is a short-lift
ratio and is not the coefficient in
\eqref{cmt-exact}. Its Ward defects and folded tails
are nonzero at finite period. This argument does not
prove exact finite marginal preservation or set the
finite classical leakage to zero. The finite reference
telescope and finite contraction remain valid without
that identity.

Equation \eqref{cmt-drift-bound} does not compute the
determinant/pivot coefficient, prove its sign or convergence
with \(M\), or control an infinite physical coupling
trajectory. The result concerns the selected joint
infinite classical system and does not establish
subsequence independence, global sector matching,
observable convergence or continuum existence.
